\begin{proposition}[C.3.23 (Proved): an additive integer before reduction modulo 28]
Fix the linear semifree circle action of C.3.14, with its
circle direction fixed. On the group of orientation-preserving
equivariant diffeomorphisms of $S^6$ there is an integer
\[
 \mathfrak q:\pi_0\operatorname{Diff}^+_{S^1}(S^6)
                    \longrightarrow\Z
\]
defined by the unframed Haefliger coordinate of the orbit
three-knot of equivariant clutching. It is a group homomorphism.
For a representative normalized to be the identity on a fixed
tubular neighbourhood, it is Wall's unframed knot coordinate,
so $\mathfrak q=-J$ in the convention of C.3.14. Its image in
the ordinary group $\Theta_7$ is the fixed signed reduction
modulo $28$ described there. In particular, if an equivariant
representative has a sixth power equivariantly isotopic to a
representative with $J=0$ as an integer, its ordinary mapping
class is zero.
\end{proposition}
\begin{proof}
We first make the normalization and its independence explicit.
The fixed set is $F=S^2$, and its equivariant normal bundle is
the trivial scalar complex rank-two bundle after conjugating
the negative weight, by C.3.14. An equivariant real derivative
on that normal representation is complex linear and has
positive real determinant. An orientation-preserving ambient
diffeomorphism therefore restricts to an orientation-preserving
diffeomorphism of $F$.

By \citet[Theorem A, page 621]{Smale1959E}, this restriction is
isotopic to the identity. Lift such an isotopy using the product
normal bundle and an invariant tubular cutoff, extending its
velocity field equivariantly. After this adjustment the map
fixes $F$ pointwise. Its normal derivative is a map
$F\to GL_2(\C)$. Polar decomposition and $\pi_2(U(2))=0$
give a homotopy to the identity, again realized in the fixed
tube by an equivariant isotopy. Finally the usual rescaling
of a tubular germ to its derivative, followed by an invariant
cutoff, makes the whole map the identity on a smaller fixed
tube. Equivariance rules out a derivative from the nontrivial
normal representation to the fixed tangent representation;
the rescaled germ is smooth at the zero parameter. This
normalization uses equivariant isotopies throughout.

Clutch two standard linear equivariant seven-discs by this
diffeomorphism. Their fixed three-discs clutch to $S^3$.
Their orbit six-discs clutch to a homotopy six-sphere, which
is standard because $\Theta_6=0$. The fixed image is thus a
smooth oriented three-knot in $S^6$. Its Haefliger integer is
the proposed $\mathfrak q$, using the connected-sum group of
oriented three-knots in the six-sphere
\citep[Sections 1 and 5.16]{Haefliger1966D}.

This integer does not depend on the chosen normalization.
An equivariant isotopy of boundary clutching maps extends
over a collar to a diffeomorphism of the clutched equivariant
pairs. The resulting orbit pairs are diffeomorphic. One can
choose compatible unitary normal tubes along the entire
isotopy, so that their Hopf orbit charts form a smooth family;
this does not invoke functorial smoothness for arbitrary
equivariant maps in a fixed raw quotient atlas.
There is also no ambiguity from identifying the orbit homotopy
sphere with $S^6$. An orientation-preserving diffeomorphism of
$S^6$ is isotopic to one supported in a ball disjoint from a
given compact three-knot: make its germ the identity near a
point, use stereographic coordinates on the complement, and
conjugate its compact support into a small ball by translations
and positive dilations, all isotopic to the identity. Such a
representative fixes that knot. Consequently the ambient
identification cannot change its knot integer.

We prove additivity geometrically. Normalize two equivariant
diffeomorphisms to be the identity near $F$. In particular each
is the identity near a chosen fixed point. Stereographic
projection from that point identifies the complement with
$\R^2\oplus\C^2$, with the circle acting scalarly on the
second summand. The diffeomorphism has compact support there.
Positive dilations and translations along the fixed $\R^2$
commute with the circle and are equivariantly isotopic to the
identity on the one-point compactification. By conjugation,
place the two supports inside disjoint small invariant balls
centered at distinct fixed points. The resulting maps remain
in their original equivariant components.

For disjoint supports, the composition clutch is the
equivariant connected sum of the two separate clutches.
This follows directly by cutting the equatorial sphere between
the two invariant balls and inserting a product collar; the
map on the separating collar is the identity. The fixed
three-discs undergo the same boundary connected-sum operation.
On the orbit pairs this is the ordinary connected sum of
the two oriented three-knots. The defining group operation
of their Haefliger integers is connected sum. Hence
\[
 \mathfrak q(\phi_1\phi_2)
 =\mathfrak q(\phi_1)+\mathfrak q(\phi_2).
\]
This proves additivity without passing first to $\Theta_7$.

For a fixed-tube-relative representative, the surgery cap and
Wall calculation of C.3.14 apply verbatim: the inherited product
normal framing gives $\mathfrak q=-J$. Altering this framing
changes only Wall's framing coordinate, whose equivariant lift
extends over the seven-disc by the explicit quaternionic
formula in C.3.14. Finally the orbit-knot generator comparison
of \citet[Example 7.14 and Proposition 7.15]{Davis1982D}
identifies the underlying ordinary sphere class with the fixed
signed reduction of $\mathfrak q$ modulo $28$.

If $\phi^6$ is equivariantly isotopic to a representative with
integer $J=0$, then $6\mathfrak q(\phi)=0$ in $\Z$.
Thus $\mathfrak q(\phi)=0$, and its underlying ordinary
mapping class is zero. This last implication uses a torsion-free
integer before reduction; it is not inferred from the weaker
equation $6[\phi]=0$ in $\Z/28$.
\end{proof}
